% Auto-generated from book/01_quaternions.md — do not edit by hand
% Regenerated by scripts/md_to_latex.py

\chapter{Chapter 1 --- Quaternions: Algebra, Geometry, and Arithmetic}
\label{ch:ch01_quaternions}

This chapter builds the algebraic and geometric language used everywhere later in the book. Chapter 2 fibers the unit sphere of quaternions; Chapters 3--4 act on it by left and right multiplication; later chapters repeatedly invoke norms, integer orders, and the four-square theorem.

\textbf{Learning goals}

\begin{enumerate}
\item Master the algebraic rules and geometric interpretation of quaternions.  
\item See unit quaternions as the three-sphere \(S^3\).  
\item Distinguish Lipschitz and Hurwitz orders and understand why the latter is preferred for arithmetic.  
\item Re-encounter the four-square theorem as a consequence of norm multiplicativity.  
\item Meet the double-cover relationship to three-dimensional rotations.  
\item Run first computational labs with \texttt{kingdom.core.quaternion}.
\end{enumerate}

\textbf{Figures in this chapter}

\begin{center}
\small
\begin{tabular}{@{}lp{0.28\textwidth}p{0.28\textwidth}@{}}
\toprule
Tag & File & Role \\
\midrule
Fig. 1.1 & \texttt{figures/fig1\_1\_ijk\_multiplication.png} & \(i,j,k\) cycle and multiplication table \\
Fig. 1.2 & \texttt{figures/fig1\_2\_s3\_unit\_quaternions.png} & Unit quaternions as \(S^3\subset\mathbb{R}^4\) \\
Fig. 1.3 & \texttt{figures/fig1\_3\_lipschitz\_hurwitz.png} & Lipschitz vs Hurwitz lattice (schematic) \\
Fig. 1.4 & \texttt{figures/fig1\_4\_norm\_multiplicativity.png} & Norm multiplicativity \(\rightarrow\) four squares \\
Aux A1.1 & \texttt{figures/aux1\_1\_quaternion\_rotation.png} & Rotation via \(v\mapsto qvq^{-1}\) (\(\rightarrow\) Ch. 2) \\
\bottomrule
\end{tabular}
\end{center}


\textbf{Claim discipline}

\begin{center}
\small
\begin{tabular}{@{}lp{0.55\textwidth}@{}}
\toprule
Claim & Type \\
\midrule
Algebra of \(\mathbb{H}\); \(N(q_1 q_2)=N(q_1)N(q_2)\); four-square theorem; \(\mathrm{Spin}(3)\to SO(3)\) double cover & \textbf{Theorem} (classical) \\
Preferring the Hurwitz order as the primary integer lattice for later flux/lattice work & \textbf{Model} choice \\
Kingdom Come / \texttt{flux\_hopf\_lib} APIs and demos & \textbf{Software fact} \\
\bottomrule
\end{tabular}
\end{center}


\bigskip
\noindent\rule{\textwidth}{0.4pt}
\bigskip

\section{1.1 Definition and basic operations}
\label{ch:ch01_quaternions:1-1-definition-and-basic-operations}

\subsection{The algebra \(\mathbb{H}\)}
\label{ch:ch01_quaternions:the-algebra-mathbb-h}

A \textbf{quaternion} is an expression

\[
q = a + bi + cj + dk,
\]

where \(a,b,c,d\in\mathbb{R}\) and the symbols \(i,j,k\) satisfy Hamilton's rules

\[
i^2 = j^2 = k^2 = ijk = -1.
\]

These force the cyclic products

\[
ij = k,\quad jk = i,\quad ki = j
\]

and the anticyclic products

\[
ji = -k,\quad kj = -i,\quad ik = -j.
\]

In particular, multiplication is \textbf{associative} and \textbf{distributive} over addition, but \textbf{not commutative}: \(ij\neq ji\).

As a real vector space,

\[
\mathbb{H} \;\cong\; \mathbb{R}^4,
\]

with basis \(\{1,i,j,k\}\). Addition and scalar multiplication are componentwise. The interesting structure is the bilinear product determined by the table above.

\begin{figure}[htbp]
  \centering
  \includegraphics[width=0.92\textwidth,height=0.42\textheight,keepaspectratio]{fig1_1_ijk_multiplication.png}
  \caption{Figure 1.1 --- Quaternion imaginary units and multiplication table.}
  \label{fig:fig1-1-ijk-multiplication}
\end{figure}

\begin{quote}\small\textit{Figure 1.1.* Left: the positive cycle \(i\to j\to k\to i\). Right: the \(4\times 4\) multiplication table for \(\{1,i,j,k\}\). Red-tinted cells are negative results; green-tinted cells are the pure imaginary basis elements.}\end{quote}

\textbf{Notation.} We write \(q = (a,b,c,d)\) or \(q = a + \mathbf{v}\) with vector part \(\mathbf{v}=(b,c,d)\in\mathbb{R}^3\). Kingdom Come stores components as \texttt{(w, x, y, z)} with \(w=a\) the real part---the same convention as \texttt{flux\_hopf\_lib.quaternion.core.Quaternion}.

\subsection{Explicit product formula}
\label{ch:ch01_quaternions:explicit-product-formula}

If \(q_1 = a_1 + b_1 i + c_1 j + d_1 k\) and \(q_2 = a_2 + b_2 i + c_2 j + d_2 k\), then

\begin{align*}
q_1 q_2
&=
(a_1 a_2 - b_1 b_2 - c_1 c_2 - d_1 d_2)\\
&\quad+ (a_1 b_2 + b_1 a_2 + c_1 d_2 - d_1 c_2)\, i\\
&\quad+ (a_1 c_2 - b_1 d_2 + c_1 a_2 + d_1 b_2)\, j\\
&\quad+ (a_1 d_2 + b_1 c_2 - c_1 b_2 + d_1 a_2)\, k.
\end{align*}

This is exactly the formula implemented by \texttt{Quaternion.multiply} in both \texttt{flux\_hopf\_lib} and \texttt{kingdom.core.quaternion}.

\subsection{Complex pairs}
\label{ch:ch01_quaternions:complex-pairs}

Identifying \(\mathbb{C}\) with \(\mathrm{span}_{\mathbb{R}}\{1,i\}\) inside \(\mathbb{H}\), every quaternion can be written uniquely as

\[
q = z_1 + z_2\, j,\qquad z_1,z_2\in\mathbb{C},
\]

or as a pair \((z_1,z_2)\in\mathbb{C}^2\). This view is the natural bridge to the Hopf map in Chapter 2.

\bigskip
\noindent\rule{\textwidth}{0.4pt}
\bigskip

\section{1.2 Conjugation and the multiplicative norm}
\label{ch:ch01_quaternions:1-2-conjugation-and-the-multiplicative-norm}

\subsection{Conjugate}
\label{ch:ch01_quaternions:conjugate}

The \textbf{conjugate} of \(q = a + bi + cj + dk\) is

\[
\overline{q} = a - bi - cj - dk.
\]

Properties (all \textbf{Theorems}, elementary verifications):

\[
\overline{\overline{q}} = q,\qquad
\overline{q_1 + q_2} = \overline{q_1}+\overline{q_2},\qquad
\overline{q_1 q_2} = \overline{q_2}\,\overline{q_1}.
\]

Note the \textbf{order reversal} under conjugation of products---the noncommutative analogue of the complex rule \(\overline{zw}=\overline{z}\,\overline{w}\).

\subsection{Norm}
\label{ch:ch01_quaternions:norm}

The \textbf{(Euclidean) norm} is

\[
N(q) \;:=\; q\,\overline{q}
\;=\; a^2 + b^2 + c^2 + d^2
\;=\; |q|^2
\]

if one writes \(|q|\) for the Euclidean length in \(\mathbb{R}^4\). In code, \texttt{Quaternion.norm()} returns \(\sqrt{N(q)}=|q|\); the multiplicative quantity used in number theory is \(N(q)=|q|^2\). We will say \textbf{multiplicative norm} for \(N(q)\) and \textbf{length} for \(|q|\) when both appear.

\textbf{Theorem (multiplicativity of the norm).} For all \(q_1,q_2\in\mathbb{H}\),

\[
N(q_1 q_2) = N(q_1)\, N(q_2).
\]

\emph{Sketch.} \(N(q_1 q_2)=(q_1 q_2)\overline{(q_1 q_2)}=q_1 q_2 \overline{q_2}\,\overline{q_1}=q_1 N(q_2)\overline{q_1}=N(q_1)N(q_2)\), using that \(N(q_2)\) is real (hence central).

\begin{figure}[htbp]
  \centering
  \includegraphics[width=0.92\textwidth,height=0.42\textheight,keepaspectratio]{fig1_4_norm_multiplicativity.png}
  \caption{Figure 1.4 --- Multiplicativity of the quaternion norm.}
  \label{fig:fig1-4-norm-multiplicativity}
\end{figure}

\begin{quote}\small\textit{Figure 1.4.* The product of two quaternions multiplies their norms. Every identity among sums of four squares is, at bottom, this theorem in arithmetic clothing.}\end{quote}

\subsection{Inverse}
\label{ch:ch01_quaternions:inverse}

If \(q\neq 0\), then

\[
q^{-1} = \frac{\overline{q}}{N(q)},
\]

so every nonzero quaternion is a unit in the division ring \(\mathbb{H}\). In floating-point code, \texttt{inverse()} divides by \(|q|^2\) and raises if \(|q|\) is near zero.

\bigskip
\noindent\rule{\textwidth}{0.4pt}
\bigskip

\section{1.3 Unit quaternions and the identification \(S^3\subset\mathbb{H}\cong\mathbb{R}^4\)}
\label{ch:ch01_quaternions:1-3-unit-quaternions-and-the-identification-s-3-subset-mathb}

\subsection{The three-sphere}
\label{ch:ch01_quaternions:the-three-sphere}

The set of \textbf{unit quaternions} is

\[
\mathrm{Sp}(1)
\;:=\;
\{ q\in\mathbb{H} : N(q)=1 \}
\;=\;
S^3 \subset \mathbb{R}^4.
\]

It is a compact Lie group under quaternion multiplication (associative, identity \(1\), inverses \(\overline{q}\)).

Geometrically:

\begin{itemize}
\item As a manifold, \(S^3\) is the unit sphere in four-dimensional Euclidean space.  
\item As a group, left (or right) multiplication by a fixed unit quaternion is an isometry of \(S^3\).  
\item As a complex pair, unit length means \(|z_1|^2+|z_2|^2=1\).
\end{itemize}

\begin{figure}[htbp]
  \centering
  \includegraphics[width=0.92\textwidth,height=0.42\textheight,keepaspectratio]{fig1_2_s3_unit_quaternions.png}
  \caption{Figure 1.2 --- Unit quaternions as \(S^3\subset\mathbb{R}^4\).}
  \label{fig:fig1-2-s3-unit-quaternions}
\end{figure}

\begin{quote}\small\textit{Figure 1.2.* Left: schematic of the constraint \(|z_1|^2+|z_2|^2=1\) via nested level sets of \(|z_1|\). Right: after stereographic projection \(S^3\setminus\{\mathrm{pt}\}\to\mathbb{R}^3\), great-circle fibers will appear as linked circles (Chapter 2; compare Figs. 0.1--0.2).}\end{quote}

\subsection{Stereographic projection (preview)}
\label{ch:ch01_quaternions:stereographic-projection-preview}

Removing a point (classically \(-1\), or any fixed unit quaternion) and projecting to a tangent \(\mathbb{R}^3\) yields a global chart on \(S^3\setminus\{\mathrm{pt}\}\). Kingdom Come's Hopf visualizer uses stereographic (or orthographic 2D) projections heavily; here we only need that **\(S^3\) is drawable** once a chart is chosen. The full Hopf fibration story---fibers, linking, base \(S^2\)---is Chapter 2.

\subsection{Normalization in software}
\label{ch:ch01_quaternions:normalization-in-software}

\texttt{Quaternion.normalize()} returns \(q/|q|\) (or the identity if \(q\) is numerically zero). Portal methods that construct new quaternions return \texttt{type(self)} so subclass helpers such as \texttt{hopf\_image()} survive \texttt{normalize()} and \texttt{multiply()}.

\bigskip
\noindent\rule{\textwidth}{0.4pt}
\bigskip

\section{1.4 Integer quaternions: Lipschitz order vs Hurwitz order}
\label{ch:ch01_quaternions:1-4-integer-quaternions-lipschitz-order-vs-hurwitz-order}

\subsection{Lipschitz order}
\label{ch:ch01_quaternions:lipschitz-order}

The \textbf{Lipschitz order} is the set of quaternions with integer coordinates:

\[
L \;:=\; \mathbb{Z}[i,j,k]
\;=\;
\{ a+bi+cj+dk : a,b,c,d\in\mathbb{Z} \}.
\]

It is a free \(\mathbb{Z}\)-module of rank \(4\), closed under multiplication, and contains the eight \textbf{Lipschitz units}

\[
\{\pm 1,\pm i,\pm j,\pm k\}.
\]

Norms of Lipschitz quaternions are nonnegative integers. Multiplicativity still holds, so products of sums of four integer squares are again such sums.

\subsection{Hurwitz order}
\label{ch:ch01_quaternions:hurwitz-order}

The \textbf{Hurwitz order} \(\mathcal{H}\) consists of all quaternions whose coordinates are \textbf{either all integers or all half-integers}:

\[
\mathcal{H}
=
L \;\cup\;
\Bigl(\tfrac12 + \tfrac12 i + \tfrac12 j + \tfrac12 k + L\Bigr).
\]

Equivalently: \(a,b,c,d\in\mathbb{Z}\) or \(a,b,c,d\in\mathbb{Z}+\tfrac12\). This is still a ring (in fact a maximal order in the rational quaternion algebra \(\mathbb{H}_{\mathbb{Q}}\)), and it is \textbf{larger} than \(L\).

The unit group of \(\mathcal{H}\) has \textbf{24} elements---the binary tetrahedral group---including

\[
\pm 1,\;\pm i,\;\pm j,\;\pm k,\quad
\tfrac12(\pm 1\pm i\pm j\pm k)
\]

(all sign combinations). The extra units and lattice points make \(\mathcal{H}\) \textbf{Euclidean} with respect to the norm: for any \(q,d\in\mathcal{H}\) with \(d\neq 0\) there exist \(m,r\in\mathcal{H}\) with \(q = md + r\) and \(N(r)<N(d)\). That Euclidean property underwrites unique factorization phenomena used in classical proofs of the four-square theorem and in ideal theory (Chapter 9).

\begin{figure}[htbp]
  \centering
  \includegraphics[width=0.92\textwidth,height=0.42\textheight,keepaspectratio]{fig1_3_lipschitz_hurwitz.png}
  \caption{Figure 1.3 --- Lipschitz vs Hurwitz (schematic slices).}
  \label{fig:fig1-3-lipschitz-hurwitz}
\end{figure}

\begin{quote}\small\textit{Figure 1.3.* Left: integer lattice (Lipschitz), shown in a 2D schematic slice. Right: Hurwitz adds half-integer points (orange squares) when all four coordinates share the same half-integer type. The picture is only a slice---true lattice geometry lives in \(\mathbb{R}^4\).}\end{quote}

\subsection{Why prefer Hurwitz later? (\textbf{Model})}
\label{ch:ch01_quaternions:why-prefer-hurwitz-later-model}

\begin{center}
\small
\begin{tabular}{@{}lp{0.28\textwidth}p{0.28\textwidth}@{}}
\toprule
Feature & Lipschitz \(L\) & Hurwitz \(\mathcal{H}\) \\
\midrule
Definition & Integer coords & Integer or half-integer coords \\
Units & 8 & 24 \\
Euclidean algorithm & fails in general & holds \\
Unique factorization flavor & weaker & stronger (classical) \\
Relation to four squares & works, clumsier proofs & clean norm arguments \\
\bottomrule
\end{tabular}
\end{center}


\textbf{Claim type: Model.} For the gauged Hopf lattice and flywheel arithmetic of Parts II--IV, this book \textbf{defaults to Hurwitz} (or to lattices containing \(\mathcal{H}\)) unless a section explicitly restricts to Lipschitz. That is a design choice aligned with classical quaternion number theory, not a theorem that Lipschitz is ``wrong.''

\textbf{Open thread.} Discrete adjacency rules built on \(\mathcal{H}\) (or on unit spheres over \(\mathcal{H}\)) feed Open Problem 1 in Chapter 3.

\bigskip
\noindent\rule{\textwidth}{0.4pt}
\bigskip

\section{1.5 The four-square theorem via norms}
\label{ch:ch01_quaternions:1-5-the-four-square-theorem-via-norms}

\textbf{Theorem (Lagrange).} Every natural number \(n\) is a sum of four integer squares:

\[
n = a^2 + b^2 + c^2 + d^2
\quad\text{for some }a,b,c,d\in\mathbb{Z}.
\]

\textbf{Quaternionic reading.} The statement is exactly: for every \(n\in\mathbb{N}\) there exists \(q\in L\) (Lipschitz) with \(N(q)=n\).

\textbf{Why multiplicativity matters.} Euler's four-square identity is \(N(q_1 q_2)=N(q_1)N(q_2)\) written in coordinates. Therefore it suffices to prove the theorem for primes (or for \(n\) square-free) and multiply representations. Classical proofs then show every prime is a norm from \(\mathcal{H}\) (or handle \(2\) and primes \(1\bmod 4\) / \(3\bmod 4\) by standard casework). Full proofs appear in many number-theory texts; Appendix D will collect a sketch. For this chapter, the essential message is:

> Four squares are not an isolated curiosity---they are the \textbf{norm theory of integer quaternions}.

\textbf{Comparison with Hatcher.} Hatcher's early chapters revolve around \textbf{two} squares and binary quadratic forms (Gaussian integers, topographs, class groups). Quaternions are the four-dimensional upgrade: the same geometric instinct (norms, units, lattices), one dimension higher. Chapter 5 will lift topographs themselves.

\textbf{Claim type: Theorem} (classical). No new proof is claimed here.

\bigskip
\noindent\rule{\textwidth}{0.4pt}
\bigskip

\section{1.6 Quaternions as rotations: the double cover \(\mathrm{Spin}(3)\to SO(3)\)}
\label{ch:ch01_quaternions:1-6-quaternions-as-rotations-the-double-cover-mathrm-spin-3-}

\subsection{Conjugation action on pure imaginaries}
\label{ch:ch01_quaternions:conjugation-action-on-pure-imaginaries}

Identify \(\mathbb{R}^3\) with the pure imaginary quaternions

\[
\mathrm{Im}\,\mathbb{H} = \{ bi + cj + dk : b,c,d\in\mathbb{R} \}.
\]

For a \textbf{unit} quaternion \(q\in S^3\) and \(v\in\mathrm{Im}\,\mathbb{H}\), set

\[
\rho_q(v) \;:=\; q\, v\, q^{-1} = q\, v\, \overline{q}.
\]

Then \(\rho_q(v)\) is again pure imaginary, and \(\rho_q\) is an orthogonal linear map of \(\mathbb{R}^3\) with determinant \(+1\). Thus we obtain a homomorphism

\[
\rho: S^3 \longrightarrow SO(3),\qquad q\mapsto \rho_q.
\]

\subsection{Axis--angle formula}
\label{ch:ch01_quaternions:axis-angle-formula}

If \(q = \cos(\theta/2) + \sin(\theta/2)\, u\) with unit pure imaginary \(u\) (i.e. \(u^2=-1\)), then \(\rho_q\) is rotation by angle \(\theta\) about the axis \(u\). This is the content of Rodrigues' formula; Kingdom Come exposes both \texttt{Quaternion.from\_axis\_angle(axis, theta)} and \texttt{rodrigues\_rotation}.

\begin{figure}[htbp]
  \centering
  \includegraphics[width=0.92\textwidth,height=0.42\textheight,keepaspectratio]{aux1_1_quaternion_rotation.png}
  \caption{Auxiliary Figure A1.1 --- Unit quaternion rotation \(v\mapsto qvq^{-1}\).}
  \label{fig:aux1-1-quaternion-rotation}
\end{figure}

\begin{quote}\small\textit{Auxiliary Figure A1.1.* A vector \(v\) swept around an axis \(u\) under the conjugation action of a one-parameter family of unit quaternions. The same rotation arises from \(q\) and from \(-q\).}\end{quote}

\subsection{The kernel is \(\{\pm 1\}\)}
\label{ch:ch01_quaternions:the-kernel-is-pm-1}

Clearly \(\rho_{-q}=\rho_q\), since \((-q)v(-q)^{-1}=qvq^{-1}\). In fact

\[
\ker\rho = \{\pm 1\},
\]

so \(\rho\) is a \textbf{2-to-1} covering homomorphism

\[
S^3 \;\cong\; \mathrm{Spin}(3) \;\longrightarrow\; SO(3).
\]

Topologically this is the unique nontrivial double cover of \(SO(3)\). Paths in \(SO(3)\) that rotate by \(2\pi\) lift to paths in \(S^3\) from \(q\) to \(-q\); only a \(4\pi\) rotation closes in the cover---the origin of spinorial sign changes in physics.

\textbf{Claim type: Theorem} (classical Lie theory / geometry).

\textbf{Forward link.} The same unit quaternions that rotate \(\mathbb{R}^3\) will, in Chapter 2, be points of the \textbf{total space} of the Hopf fibration. Left and right multiplications that preserve the fibration become the symmetry language of Chapters 3--4.

\bigskip
\noindent\rule{\textwidth}{0.4pt}
\bigskip

\section{1.7 First computational examples and labs}
\label{ch:ch01_quaternions:1-7-first-computational-examples-and-labs}

Core API: \texttt{kingdom.core.quaternion} (see also Appendix C ?C.1).

\textbf{Labs (short form).} Full listings: \textbf{Appendix C}.

\begin{itemize}
\item \textbf{1.A} Verify \(ij=k\), \(ji=-k\), \(i^2=-1\) with \texttt{Quaternion}.
\item \textbf{1.B} Check \(N(q_1q_2)=N(q_1)N(q_2)\) numerically (residual \(\sim 10^{-14}\)).
\item \textbf{1.C} Enumerate small Lipschitz norms.
\item \textbf{1.D} Confirm \(q\) and \(-q\) induce the same conjugation action on a pure vector.
\item \textbf{1.E} Optional: \texttt{from\_hopf\_coords} / \texttt{hopf\_image} (preview of Ch. 2).
\end{itemize}

\begin{lstlisting}[style=qga,language=python]
from kingdom.core.quaternion import Quaternion
i, j = Quaternion(0,1,0,0), Quaternion(0,0,1,0)
print(i.multiply(j))  # k
\end{lstlisting}

\bigskip
\noindent\rule{\textwidth}{0.4pt}
\bigskip

\section{Exercises}
\label{ch:ch01_quaternions:exercises}

\textbf{1.A (hand).} Expand \((a+bi+cj+dk)(e+fi+gj+hk)\) and match the formula in ?1.1.

\textbf{1.B (hand).} Prove \(\overline{q_1 q_2}=\overline{q_2}\,\overline{q_1}\) from the multiplication table (or from the explicit formula).

\textbf{1.C (hand).} Show that \(N(q)=q\overline{q}\) is real and nonnegative, and that \(N(q)=0\Leftrightarrow q=0\).

\textbf{1.D (hand).} Verify that the 24 Hurwitz units listed in ?1.4 all have norm \(1\). How many Lipschitz units are among them?

\textbf{1.E (hand / reading).} Using multiplicativity, show that if \(n=N(q_1)\) and \(m=N(q_2)\) then \(nm=N(q_1 q_2)\). Explain in one sentence why this reduces Lagrange's theorem to representing primes.

\textbf{1.F (hand).} Let \(q=\cos(\theta/2)+\sin(\theta/2)\,i\). Compute \(q\, j\, \overline{q}\) and identify the rotation in the \(jk\)-plane.

\textbf{1.G (code).} Complete Labs 1.A--1.D. Submit residuals for Lab 1.B and a printed set of norms for Lab 1.C with \(N_{\max}=3\).

\textbf{1.H (code).} Using \texttt{chordal\_distance} (upstream) or \(\min(|q-r|,|q+r|)\), verify numerically that \(q\) and \(-q\) are close in the rotation sense even when far in \(\mathbb{R}^4\).

\textbf{1.I (Hatcher bridge).} In Hatcher, sums of two squares and Gaussian integers organize binary quadratic forms. Write three sentences mapping: (Gaussian integers \(\leftrightarrow\) Lipschitz/Hurwitz), (complex modulus \(\leftrightarrow\) quaternion norm), (unit circle \(S^1\) \(\leftrightarrow\) unit quaternions \(S^3\)).

\textbf{1.J (forward).} Without reading Chapter 2, guess: if unit quaternions are points of \(S^3\), what might a natural map \(S^3\to S^2\) do to a whole circle of phases? (You will check your guess against the Hopf map next.)

\bigskip
\noindent\rule{\textwidth}{0.4pt}
\bigskip

\section{Code and asset pointers}
\label{ch:ch01_quaternions:code-and-asset-pointers}

\begin{lstlisting}[style=qga]
kingdom.core.quaternion.Quaternion
kingdom.core.quaternion.rodrigues_rotation
flux_hopf_lib.quaternion.core.Quaternion   # upstream implementation
flux_hopf_lib.quaternion.core.q_mult, q_conj, q_normalize
\end{lstlisting}

\textbf{Figures:} generated for this chapter under \texttt{book/figures/fig1\_*} and \texttt{aux1\_1\_*} (see \texttt{FIGURES.md}). \textbf{Related portal:} Gradio rotation intuition is secondary here; Hopf Visualizer becomes primary in Chapter 2.

\bigskip
\noindent\rule{\textwidth}{0.4pt}
\bigskip

\section{Looking ahead}
\label{ch:ch01_quaternions:looking-ahead}

With the algebraic and geometric language of quaternions in hand---product, conjugate, norm, \(S^3\), Hurwitz arithmetic, and the double cover of \(SO(3)\)---we are ready to \textbf{fiber} \(S^3\) over \(S^2\) in Chapter 2. The same unit quaternions will reappear as total-space points; their left and right multiplications will become the symmetries of the gauged Hopf lattice in Part II.

\bigskip
\noindent\rule{\textwidth}{0.4pt}
\bigskip


